We study machine-learning detection of controlled beyond-General-Relativity
(beyond-GR) deviations in gravitational-wave signals, using both synthetic
aLIGO-PSD noise and real LIGO H1 detector strain. Three deviation families
are applied to General-Relativistic inspiral-merger-ringdown waveforms:
amplitude modulation, phase modulation, and frequency modulation, each
parameterized by a dimensionless strength coefficient $\beta$. A hybrid
classifier combining a one-dimensional convolutional neural network with
ten hand-crafted waveform statistics is trained on GR and modified
waveforms and tested on a deviation type excluded from training. The
central result is a quantitative detectability curve as a function of
$\beta$. Using the real GW150914 strain as a template and real H1
detector noise, we find a detection threshold at $\beta \approx 0.25$,
with accuracy rising smoothly from chance at $\beta \leq 0.2$ to perfect
classification at $\beta \geq 0.5$. The threshold value is specific to
the quadratic-in-time modulation form adopted here and should not be
interpreted as a generic constraint on beyond-GR parameters. We
nevertheless argue that the negative result at small $\beta$ is
informative: it establishes a quantitative limit on machine-learning-only
beyond-GR searches in real detector noise, in the absence of matched-filter
signal extraction.
